\subsection*{Chapter 36: Project 5: Reasoning RL and Alignment}

\paragraph{Exercise 36.1.} One configuration-seed is $200\times40\,\text{s}=8000$\,s on two GPUs.
Four variants $\times$ four seeds give $16\times8000\,\text{s}\times2\,\text{GPUs}\approx71$
GPU-hours. Pruning 30\% of the 14\,s training phase saves 4.2\,s per step, a step of 35.8\,s, or
about 64 GPU-hours (10\% less). Generation dominates, so pruning helps much less than its 30\%
suggests.

\paragraph{Exercise 36.2.} Per rollout batch: 26\,s of generation, then $32\times0.5=16$\,s of
training (ignoring grading and sync). The cycle is 42\,s; GPU~1 is idle for $16/42\approx38\%$ of it
and GPU~0 for $26/42\approx62\%$. Optimizer steps per hour rise from $3600/40=90$ to
$32\times3600/42\approx2700$, a $30\times$ increase in update frequency for the same rollouts. Whether
this helps depends on how stale updates behave; that is what the experiment measures.

\paragraph{Exercise 36.3.} With independent misses at rate $m=0.04$, an observed success occurs with
probability $0.5\times0.96=0.48$. Correct-but-missed responses receive the advantage of a failure
(about $-0.48$ instead of $+0.52$), so their tokens are pushed down. On average this scales the
useful signal and adds noise, but it does not bias the direction. If misses concentrate in long
responses, the noise becomes a systematic penalty on long correct reasoning: the policy learns to
shorten or reformat answers for the parser rather than to be right. Inspect misses by length before
trusting length trends.

\paragraph{Exercise 36.4.} Colocated: $26+14+5=45$\,s per step on one GPU, versus about 40\,s on two
GPUs. Wall-clock is 12.5\% slower, but GPU-seconds per step drop from $80$ to $45$, so colocation is
about $1.8\times$ more efficient per GPU-hour whenever the phases cannot overlap anyway.

\paragraph{Exercise 36.5.} Policy GPU: bf16 weights 16\,GB, bf16 gradients 16\,GB, RMSprop square
averages in fp32 32\,GB, and optionally fp32 master weights 32\,GB, for 64--96\,GB. Activations for
one pair of 512-token sequences are a few GB with FlashAttention, and the logits are about
$2\times512\times128{,}256\times4\,\text{B}\approx0.5$\,GB in fp32. The reference GPU needs about 16\,GB of
weights plus small activations. To fit in 80\,GB, first drop the fp32 master copy (or keep RMSprop
state in bf16), then use activation checkpointing. Sharding is the next step if that is not enough.
